\section{本讲小结}
角动量同时依赖动量和参考点，守恒依据是相应外力矩或角冲量为零。固定支点可破坏线动量守恒而保留绕轴角动量；中心力使面积速度不变；绕竖直轴守恒的分量不能扩大成任意轴都守恒。

综合问题先确定系统、参考点和事件，再让线动量、角动量、能量分别提供适当的方程。相对速度要转换为同一惯性系速度，末态还需满足分离、接触和运动顺序。若优化起跳速率或求再次相遇时间，必须写出相对位移及相对加速度。
